\documentclass[12pt]{article}
\usepackage{amsmath,amsfonts,fullpage,amsthm,graphicx}

\newtheorem*{theorem}{Theorem}
\newtheorem*{fact}{Fact}
\newtheorem*{goal}{Goal}
\newtheorem*{idea}{Idea}
\newtheorem*{defn}{Definition}
\newtheorem*{ex}{Example}

%Style section
%\setlength{\textheight}{10in}
%\setlength{\textwidth}{7in}
%\hoffset=-0.75in
 \pagestyle{plain}
% \setlength{\topmargin}{0in}
% \setlength{\headsep}{0in}
% \setlength{\oddsidemargin}{0in}
% \setlength{\evensidemargin}{0in}

\begin{document}


\pagestyle{empty} %use this to get rid of page numbers

\noindent
{Math 166 \hskip 1.5 truein  Name:\ \hrulefill}

\noindent
%\vskip 0.3truein
%\bigskip
%\noindent
{Polar coordinates}
%{ and center of mass}

\noindent  \hskip 3.5 truein  Section Number:\ \hrulefill
\medskip

\begin{goal} \rm
Apply calculus to polar coordinates.
% and center of mass.
\end{goal}
%\smallskip

\begin{enumerate}
\item  Plot the following polar coordinates:  $(-7, \frac{\pi}{2})$, $(3, \frac{4\pi}{3})$, $(2, \pi)$, $(0, 2)$. Then write
each polar coordinate in a different way (i.e. with a different value of $r$ and/or $\theta$).
Lastly, from your graph or using the equations $x = r \cos(\theta)$ and $y = r \sin (\theta)$, convert the points from polar to rectangular coordinates.
\vfill
\item Plot the following rectangular coordinates: $(-5, 0)$, $(1, -2)$, $(-1, 2)$, $(12, 5)$. 
Then, convert the points from rectangular to polar coordinates. 
\vfill
\pagebreak
\item Convert $2x^2+5x^3=1+xy$ into polar coordinates. 
\vfill
\item   Convert $r = -8\cos(\theta)$ into Cartesian/rectangular coordinates.
\vfill
\item  Plot $r = -8\cos(\theta)$ for $0\leq \theta \leq 2\pi$.
\vfill
\pagebreak
\item Graph the polar equation $r = 1+2 \cos(\theta)$.
\vfill
\item Find all values of $\theta$ such that $0\leq \theta\leq 2\pi$ and $r = 1+2 \cos(\theta) = 0$.
\vfill
\pagebreak
\item  Determine the area of the inner loop of 
$r=1+2 \cos(\theta)$.
\\
\textit{Hint: Use Problems 6 and 7 to find the bounds on the integral.}
\vfill
\item  Find the arc length of the polar equation $r = 2\theta$ for $0 \leq \theta \leq 1$.
\vfill
\end{enumerate} 

\end{document}
